\documentclass{article}
\usepackage[pdftex]{graphicx}
\usepackage{xcolor}
\usepackage[letterpaper, margin = 2.5cm]{geometry}
\usepackage[T2A]{fontenc}
\usepackage[english, russian]{babel}
\usepackage{amsmath, amsfonts, amssymb}
\usepackage{hyperref, url}

\usepackage{minted}
\usepackage{parskip}

\begin{document}

\title{Домашнее задание 2}
\date{}
\author{SMR-2025 AIM, Илья Осокин @elijahmipt} %поменяйте на своё
\maketitle

Задачу 1 нужно будет оформить в tex и прислать в виде pdf (можно в Markdown в юпитер ноутбуке, но все векторы должны быть оформлены, как в техе).
В этой задаче требуется полностью письменно ответить на вопросы, то есть если спрашивают, какой вектор, напишите формулку для вектора, в которой есть все нужные компоненты, и так далее.
Задачи 2, 3 и 4 нужно будет запрогать и прислать их в виде jupyter notebook (файла .ipynb).
Во 2 задаче требуется записать короткое видео, так что в классруме нужно будет прикрепить три файла (или два, если теория оформляется в Markdown).

\section*{Задача 1}

Выведите и линеаризуйте динамику двух маятников, связанных пружиной.
Нулевой угол у них, когда они висят вниз, пружина в этом состоянии не деформирована (и создает нулевую силу).
Жесткость ее $k$, масса точечных грузов $m$, длина $l$, они в поле тяготения $g$.

\begin{enumerate}
    \item Какой у этой системы вектор состояния? Почему? Какой он размерности?
    \item Запишите кинетическую энергию системы. Запишите потенциальную энергию системы. Для потенциальной действуйте в упрощении, пренебрегая отклонением пружины от горизонтали. Считайте, что углы малые (но пока что только в этом месте) и что растяжение пружины можно найти через длину маятника и два синуса
    \item Запишите лагранжиан системы и два уравнения движения, полученные с помощью уравнения Эйлера-Лагранжа. Выразите из них вторые производные углов
    \item Линеаризуйте систему в окрестности нулевых углов, описывая, чем и почему пренебрегаем. Запишите систему в виде $\dot{x} = A x$
\end{enumerate}

\section*{Задача 2}

Реализуйте симуляцию этой системы. Рисование можно взять из семинара, просто маятников будет два. Сделайте короткую (до 10 секунд) запись поведения системы и приложите его в классруме

\section*{Задача 3}

Реализуйте функцию, вычисляющую стоимость траектории.
На вход ей поступает лист координат от времени, лист скоростей от времени, лист сил от времени и dt. На выход она отдает стоимость траектории.
Стоимость всего квадратична, находиться в $m$ метрах от начала координат в течение одной секунды стоит $3 m^2$, двигаться со скоростью $v$ метров в секунду в течении одной секунды стоит $5 v^2$, прикладывать силу в $n$ ньютонов в течение одной текунды стоит $8 n^2$.
Протестируйте свою функцию на траектории точки с массой из семинара, которую стабилизирует PD-контроллерю

\section*{Задача 4}

Попробуйте минимизировать стоимость траектории, варьируя $K_p$ и $K_d$ (можно руками, можно с помощью grid search, можно градиентным спуском, можно еще как-то, про что я не догадался).

\end{document}
